\documentclass[review]{vgtc}

\graphicspath{{figures/}{./}}
\usepackage{amsmath}
\usepackage{amssymb}
\usepackage{booktabs}
\usepackage{multirow}
\usepackage{tabularx}
\usepackage{subcaption}
\usepackage{xcolor}
\usepackage{enumitem}
\usepackage{fontspec}
\usepackage{xeCJK}
\setmainfont{Times New Roman}
\setsansfont{Arial}
\setmonofont{Courier New}
\setCJKmainfont[BoldFont=SimHei,ItalicFont=KaiTi]{SimSun}
\setCJKsansfont{Microsoft YaHei}
\setCJKmonofont{FangSong}

\onlineid{0}
\vgtccategory{Research}
\vgtcinsertpkg

\newcommand{\system}{EgoAnchor}
\newcommand{\SE}{\mathrm{SE}(3)}

\title{EgoAnchor：将异步6DoF视觉观测转化为消费级混合现实中的稳定动态物体锚点}
\author{匿名作者}
\authorfooter{}

\abstract{%
以真实物体为中心的混合现实应用要求虚拟内容在用户转头、物体移动、短时遮挡和重新获取过程中持续附着于目标。现有6DoF视觉方法通常输出图像采集时刻的相机系位姿；其异步到达、更新频率不规则和观测质量波动，使该输出难以直接作为应用锚点。本文提出\system，一套面向消费级混合现实的观测到锚点系统，将异步相机系位姿转换为连续可查询的世界系动态物体锚点。系统通过采集时刻世界对齐恢复观测的时空语义，以可视度门控的颜色--深度一致性评分控制状态更新，并结合运动状态估计、延迟Linear/SLERP时序合成、显式静止锚定与生命周期管理，定义静止、运动、遮挡和恢复阶段的输出行为。我们在Meta Quest 3上实现\system。平台参考下的受控实验表明，\system将静止平移泄漏的中位P95降至0.82~mm、静止旋转泄漏降至0.33$^{\circ}$，动态轨迹在滞后对齐后的平移与旋转RMSE分别为9.12~mm和4.81$^{\circ}$。组件实验进一步表明，采集时刻对齐、StaticLock、可靠性接纳和历史插值分别解释了世界一致性、静止稳定性和动态轨迹质量的主要改善。在24名参与者、三类日常刚性物体的被试内研究中，\system在静止稳定、位置正确、恢复一致、依赖意愿以及稳定--响应平衡上获得更高评价，并提高了参与者对空间标注的信任。上述结果表明，动态物体锚定应被视为一项具有明确稳定性、响应性与可恢复性权衡的运行时系统问题。
}

\keywords{混合现实, 动态物体锚定, 6DoF位姿估计, 异步视觉, 运行时系统, 用户信任}

\begin{document}
\firstsection{引言}
\maketitle

透视混合现实正在将操作提示、状态信息和交互控件直接放置在用户身边的真实物体上。此类界面不仅需要识别目标，还需要一个随物体运动持续更新的对象级空间参考。当用户转动头部、拿起并旋转物体、将其重新放置，或目标短暂离开视野时，虚拟内容应保持与物体一致的空间关系。传统世界锚点适用于静态环境，却无法描述物体自身的运动；单帧位姿估计则只给出特定图像时刻下的相机相对位姿。二者之间缺失的是可由应用持续查询的\emph{动态物体锚点}。

近年来，开放词表检测、视频对象分割、双目深度重建和零样本模型配准显著降低了新物体的接入成本~\cite{posecnn2017,densefusion2019,cosypose2020,megapose2022,foundationpose2023}。给定目标三维模型，视觉后端已经能够在无需逐物体训练的条件下生成6DoF位姿候选。与此同时，平台对象追踪和商业模型目标系统展示了对象级空间参考的应用价值~\cite{appleObjectAnchor,azureObjectAnchors,vuforiaModelTargets,metaDynamicObjectTracker}。然而，前者主要关注相机坐标系中的位姿估计，后者的对象范围、初始化流程和内部滤波由平台定义。对于第三方异步视觉观测，如何形成稳定、连续且具有明确失效语义的应用锚点，仍缺少统一的系统设计。

这一差距集中体现为三个相互关联的问题。首先，视觉结果在时刻$t_a$到达应用，但对应的是图像采集时刻$t_f$的相机坐标系；若直接使用到达时刻的头显位姿完成世界坐标复合，用户头动会被写入物体轨迹。其次，遮挡、深度伪影和模型错配可能产生数值形式合法、但不适合更新锚点的候选。最后，低频且不规则的观测无法直接驱动高帧率渲染，运行时需要定义观测之间、运动起停、短时丢失和重新获取时的输出行为。因此，视觉位姿观测并不天然构成可用的混合现实锚点。

本文提出\system，一套运行于消费级头显的观测到锚点系统。感知后端从双目透视图像与目标三维模型中生成相机系位姿和可靠性诊断；头显端运行时依据帧标识恢复采集时刻的相机世界变换，在统一边界上筛选观测并更新运动状态，随后通过历史时序合成生成逐渲染帧轨迹。对于静止物体，系统以显式状态抑制残余抖动；对于遮挡和恢复，生命周期管理规定继续输出、冻结、不确定标记和重新注册的条件。该架构把视觉算法视为可替换的观测生产者，将研究重点放在观测如何成为应用锚点。

我们从完整系统行为、组件归因和感知效用三个层面评价\system。受控实验使用Quest Touch Plus控制器的平台位姿作为同步参考，量化头动下的世界一致性、静止稳定性、动态轨迹质量、起动转换和遮挡行为。随后，在相同候选流和渲染时间线上逐项切换关键机制。最后，通过24名参与者和三类日常物体的被试内研究，检验这些系统差异是否影响用户对锚点质量和可信度的判断。

本文的贡献包括：
\begin{itemize}[leftmargin=*]
  \item \textbf{观测到锚点的系统架构：}将异步相机系6DoF位姿转换为连续可查询的世界系动态物体锚点，并明确感知后端与头显运行时之间的时空契约。
  \item \textbf{面向动态锚定的运行时机制：}结合采集时刻世界对齐、可靠性接纳、历史轨迹合成、显式静止锚定和生命周期管理，统一处理世界一致性、静止稳定性、运动连续性与追踪恢复。
  \item \textbf{分层系统评价：}在同一候选流下量化完整系统及关键组件的作用，并通过用户研究建立运行时行为、感知质量和空间标注信任之间的联系。
\end{itemize}

\begin{figure*}[t]
\centering
\fbox{\parbox[c][1.45in][c]{0.96\textwidth}{\centering
\textbf{Teaser图占位。} 左：用户在Quest 3透视环境中操作日常物体；中：头动、物体运动和遮挡恢复三个关键时刻，对比Arrival-Hold、One-Euro Anchor与\system的虚实叠加；右：0.82~mm静止平移泄漏、9.12~mm滞后对齐动态RMSE、18/24参与者更信任\system空间标注。}}
\caption{\system将延迟且不规则的视觉位姿观测转化为稳定的动态物体锚点。Teaser应突出“观测输入--运行时处理--应用结果”的主线。}
\label{fig:teaser}
\end{figure*}

\section{相关工作}

\subsection{混合现实中的真实物体追踪与锚定}
物理标记和物体侧追踪器能够提供稳定的目标位姿，但会改变目标外观或增加部署成本~\cite{kato1999artoolkit,olson2011apriltag,garridojurado2014aruco}。商业对象锚定系统通常结合目标模型与设备感知，在工业维护、空间标注和物体识别中提供对象级坐标系~\cite{appleObjectAnchor,azureObjectAnchors,vuforiaModelTargets,metaDynamicObjectTracker}。这些系统在支持对象、初始化方式、动态运动能力和中间观测开放程度上存在差异。\system关注的是另一类场景：应用拥有第三方视觉位姿流，需要在头显端建立统一且可解释的动态锚点语义。

\subsection{基于模型的6DoF物体位姿估计}
学习式6DoF估计已由实例特定网络逐步发展到面向未知物体的模型匹配和零样本配准~\cite{posecnn2017,densefusion2019,cosypose2020,megapose2022,foundationpose2023}。开放词表检测和视频对象分割改善了目标初始化与时序掩膜，双目重建则为模型对齐提供米制几何。上述方法的典型输出是某一图像时刻下的相机系位姿。\system将这些模块作为可替换的感知组件，并研究其下游所需的时间对齐、观测接纳、连续输出和生命周期语义。

\subsection{时延、滤波与虚实配准}
增强现实中的配准误差与感知、传输和渲染时延密切相关~\cite{azuma1994improving,diluca2010latency}。头部位姿预测和异步时间扭曲主要降低显示链路中的运动到光子时延~\cite{vanwaveren2016atw}；One-Euro等滤波方法则在含噪输入上调节抖动与响应速度。动态物体锚定还包含两项独立因素：观测属于历史相机位姿，且目标可能在采集后继续运动。采集时刻对齐恢复前者的世界坐标语义，时序合成则在轨迹误差与输出延迟之间作出选择。

\section{系统设计}

\subsection{设计目标与观测契约}
\system围绕四项系统属性展开。\emph{世界一致性}要求观测使用图像采集时刻的头显位姿进入世界坐标系；\emph{选择性更新}要求有害候选在扰动锚点前被识别；\emph{连续输出}要求低频观测之间仍能为每个渲染时刻提供确定的位姿；\emph{生命周期可预测性}要求系统在静止、运动起始、短时遮挡、持续丢失和恢复阶段采用明确策略。

对于帧标识$f$对应的双目图像$I_f=(I_f^L,I_f^R)$，感知后端返回观测
\begin{equation}
  o_f = \left(f,t_f,T_o^{c_f},R_f,d_f\right),
\end{equation}
其中$T_o^{c_f}\in\SE$表示目标在采集相机坐标系中的位姿，$R_f\in[0,1]$表示运行时可靠性分数，$d_f$包含用于日志和恢复的诊断量。头显在采集图像时记录相机世界变换$T_c^w(f)$。帧标识是两端同步的主键，候选到达时间只用于测量传输与处理延迟。

\begin{figure*}[t]
\centering
\fbox{\parbox[c][1.65in][c]{0.96\textwidth}{\centering
\textbf{Pipeline图占位。} 建议采用上下两条泳道。感知后端：双目采集$\rightarrow$目标掩膜$\rightarrow$米制深度$\rightarrow$模型配准$\rightarrow$6DoF位姿与VCD诊断。头显运行时：历史相机位姿检索$\rightarrow$采集时刻世界对齐$\rightarrow$观测接纳$\rightarrow$状态估计$\rightarrow$延迟Linear/SLERP合成$\rightarrow$StaticLock与生命周期$\rightarrow$应用锚点。图上标出$t_f$、$t_a$和$t_r$。}}
\caption{\system的观测到锚点流程。感知后端输出与采集时刻$t_f$绑定的相机系位姿；头显运行时恢复世界语义，并在渲染时刻$t_r$生成连续锚点。}
\label{fig:pipeline}
\end{figure*}

\subsection{感知后端与可靠性诊断}
默认感知后端组合开放词表目标分割、时序掩膜传播、双目深度重建和零样本模型位姿追踪~\cite{yoloe2025,cutie2024,foundationstereo2025,foundationpose2023}。这些模块为运行时提供米制6DoF候选与可用于验证的图像、深度和可视性信息。

为评估候选是否适合更新锚点，系统在估计位姿下渲染目标模型，并计算可视度门控的颜色--深度一致性分数（Visibility-gated Color-Depth, VCD）：
\begin{equation}
R_f=V_f\cdot\exp\left(\frac{w_C\log(C_f+\epsilon)+w_D\log(D_f+\epsilon)}{w_C+w_D}\right),
\end{equation}
其中$V_f$表示渲染投影得到的有效可视支持，$C_f$和$D_f$分别表示掩膜区域内的颜色一致性与米制深度一致性。几何平均使高分同时依赖外观和几何证据。VCD在本文中承担观测排序与接纳功能；通道实现、阈值和参数灵敏度放入补充材料。

\section{观测到锚点运行时}

\subsection{采集时刻世界对齐}
候选$o_f$到达后，运行时检索同一帧标识对应的相机世界变换，并计算
\begin{equation}
  T_o^w(f)=T_c^w(f)T_o^{c_f}.
  \label{eq:capture-align}
\end{equation}
在复合前，系统通过固定基变换统一OpenCV与Unity的坐标约定。式~\ref{eq:capture-align}使历史观测在其原始采集时刻进入世界坐标系，从而避免图像采集与候选到达之间的头显运动污染物体轨迹。目标在$t_f$之后发生的真实运动由后续观测和时序策略共同反映。

\subsection{观测接纳与状态估计}
只有满足VCD与当前生命周期条件的观测才进入运动状态。平移采用常速度卡尔曼模型估计位置与线速度；旋转在局部切空间中估计姿态增量与角速度。可靠性分数决定候选是否进入滤波状态，而非连续改变测量噪声。该设计保持观测筛选与状态估计的职责分离，使同一组已接纳观测对应确定的滤波行为。

\subsection{历史时序合成}
默认运行时不将每个候选外推到当前时刻，而是在渲染时刻$t_r$查询一个延迟的历史目标时刻$t_r-\Delta$。$\Delta$依据近期采集到渲染延迟自适应更新，并设置下限和变化速率约束。当目标时刻落在两个已接纳控制点之间时，平移采用线性插值，旋转采用最短弧SLERP。该策略以固定相位延迟换取有界、连续的轨迹，降低外推在速度突变和观测校正处产生的过冲。

实验中另设平滑卡尔曼外推和Hermite插值。三种策略共享相同的候选、采集时刻对齐、VCD接纳、状态估计和生命周期，仅改变渲染时刻轨迹的生成方式。

\subsection{静止锚定与生命周期}
一般滤波器在静止期抑制抖动时往往同时降低运动起始的响应速度。\system为此引入StaticLock：当线速度、角速度和可靠性在驻留时间内满足静止条件，运行时将显示锚点锁定在由近期稳定观测形成的稳健原点。持续运动、相对观测共识的位移、可靠性下降或累计方向性运动证据可触发解锁；解锁过程在冻结轨迹与运动轨迹之间平滑衔接。

生命周期状态机区分正常追踪、短时继续输出、不确定冻结、丢失和恢复。短暂缺失时，运行时利用既有历史维持输出；缺失延长后，锚点冻结于最近的可靠位置并暴露不确定状态；持续丢失触发重新注册。该状态机不仅管理内部算法，也规定应用在每种状态下获得的位姿和置信信息。

\section{系统实现}
头显前端运行于Meta Quest 3与Unity，负责双目透视图像采集、相机标定、采集时刻设备位姿记录和逐帧锚点更新。感知后端运行于配备NVIDIA RTX 3090的工作站，通过无线局域网接收图像并返回异步观测。高带宽图像流与低带宽位姿、状态和控制消息使用独立通道。头显端运行时以C\#实现，感知后端以Python实现。默认配置下，应用日志记录的候选率约为12.9~Hz，渲染端以头显帧率持续查询锚点。完整软件版本、网络消息格式和参数配置列于补充材料。

\section{系统评估}
评价分为三个层次。实验一测量应用实际接收的完整锚点行为；实验二在相同候选流和参考轨迹上切换单一设计因素；实验三在日常物体上评价感知质量与信任。实验一和实验二以预定义动作片段为统计单位，帧用于重建轨迹，不作为独立样本。

\subsection{实验一：端到端锚点行为}

\paragraph{装置与参考。}
我们使用Quest Touch Plus控制器作为视觉目标，并将平台SDK提供的控制器位姿作为同步参考。通过标定获得控制器跟踪坐标系与目标网格之间的固定变换。该参考与头显共享平台世界坐标系，适合比较同一系统内部的运行时配置；因此，以下结果表示平台参考下的相对锚点行为。

\paragraph{比较条件。}
\textbf{Arrival-Hold}在候选到达时使用当前头显位姿完成世界坐标复合，并保持到下一候选。\textbf{Capture-Hold}改用图像采集时刻的头显位姿，但不进行连续平滑。\textbf{One-Euro Anchor}采用采集时刻对齐、VCD接纳、One-Euro滤波、与完整系统相同的延迟渲染接口和生命周期，但不启用StaticLock。\textbf{EgoAnchor}采用卡尔曼状态估计、Linear/SLERP历史合成、StaticLock和完整生命周期。

\paragraph{指标。}
静止片段报告中心化平移与旋转P95，以分离固定标定偏差和头动造成的轨迹泄漏，并以相邻渲染帧增量P95衡量显示抖动。动态片段联合报告有效滞后、滞后对齐RMSE和对齐残差的帧间增量P95。有效滞后描述响应相位，滞后对齐RMSE描述移除相位差后的轨迹保真度；两者共同刻画动态锚定表现。系统还报告当前时刻RMSE、遮挡期误差和从StaticLock进入可见运动跟随的转换时间。

\begin{figure*}[t]
\centering
\begin{subfigure}[t]{0.245\textwidth}\includegraphics[width=\linewidth]{figure2a_static_translation.pdf}\caption{静止平移}\end{subfigure}\hfill
\begin{subfigure}[t]{0.245\textwidth}\includegraphics[width=\linewidth]{figure2b_static_rotation.pdf}\caption{静止旋转}\end{subfigure}\hfill
\begin{subfigure}[t]{0.245\textwidth}\includegraphics[width=\linewidth]{figure2c_dynamic_translation.pdf}\caption{动态平移}\end{subfigure}\hfill
\begin{subfigure}[t]{0.245\textwidth}\includegraphics[width=\linewidth]{figure2d_dynamic_rotation.pdf}\caption{动态旋转}\end{subfigure}
\caption{端到端锚点行为。小点表示动作片段，大标记和误差线表示中位数与四分位距。静止误差采用中心化P95；动态误差采用滞后对齐RMSE，动态抖动由对齐后的残差轨迹计算，因而不将目标真实运动计入抖动。}
\label{fig:exp1}
\end{figure*}

\begin{table*}[t]
\centering
\caption{实验一各动作片段的中位数 [Q1, Q3]。数值越低越好。}
\label{tab:exp1}
\small
\begin{tabular}{lcccccc}
\toprule
方法 & 静止平移P95 (mm) & 静止旋转P95 ($^\circ$) & 静止平移抖动 (mm) & 动态平移RMSE (mm) & 动态旋转RMSE ($^\circ$) & 动态平移残差抖动 (mm) \\
\midrule
Arrival-Hold & 43.67 [38.12, 55.80] & 6.49 [4.54, 8.32] & 11.70 [10.73, 16.76] & 18.57 [14.89, 21.66] & 6.09 [5.76, 6.51] & 41.08 [32.60, 46.05] \\
Capture-Hold & 13.03 [11.45, 16.26] & 5.40 [4.42, 7.67] & 4.30 [3.65, 6.79] & 17.14 [13.32, 19.13] & 5.84 [5.56, 6.04] & 39.69 [33.47, 45.87] \\
One-Euro & 10.65 [8.29, 13.65] & 3.29 [3.01, 4.38] & 0.85 [0.70, 1.39] & 15.69 [13.20, 20.50] & 5.55 [5.11, 6.13] & 5.25 [4.21, 7.46] \\
\system & \textbf{0.82 [0.58, 0.92]} & \textbf{0.33 [0.24, 0.36]} & \textbf{0.06 [0.06, 0.07]} & \textbf{9.12 [7.41, 11.61]} & \textbf{4.81 [4.71, 5.29]} & \textbf{4.56 [4.07, 6.14]} \\
\bottomrule
\end{tabular}
\end{table*}

\paragraph{结果。}
\system在静止平移、静止旋转和逐帧抖动上均取得最低值（表~\ref{tab:exp1}）。其静止平移泄漏的中位P95为0.82~mm，低于One-Euro的10.65~mm；静止平移帧间增量为0.06~mm。静止片段的绝对平移注册误差中位数为6.60~mm，说明较低的中心化误差同时保持在合理的绝对配准范围内。遮挡期间，\system的平移P95为4.85~mm。

在动态片段中，\system的平移与旋转有效滞后中位数分别为360~ms和345~ms，滞后对齐RMSE为9.12~mm和4.81$^{\circ}$；One-Euro的对应结果为380~ms、385~ms、15.69~mm和5.55$^{\circ}$。在不进行滞后补偿的当前时刻比较中，\system的平移与旋转RMSE为125.83~mm和31.88$^{\circ}$，One-Euro为117.61~mm和34.10$^{\circ}$。结果显示，\system主要改善的是稳定历史轨迹与旋转一致性，而当前时刻跟随仍受处理与时序合成延迟影响。StaticLock带来的起动转换时间中位数为510.36~ms，构成静止稳定性的响应代价。

\subsection{实验二：组件归因}
我们在相同观测、参考轨迹和渲染时间线上切换单一设计因素。采集时刻与到达时刻复合用于衡量时间语义；启用与关闭StaticLock用于衡量显式静止状态；VCD通过事件级风险--覆盖率曲线评价；Linear/SLERP、Hermite插值和平滑卡尔曼外推用于比较逐帧轨迹合成策略。

\begin{figure*}[t]
\centering
\begin{subfigure}[t]{0.20\textwidth}\includegraphics[width=\linewidth]{figure3a_capture_alignment.pdf}\caption{采集时刻对齐}\end{subfigure}\hfill
\begin{subfigure}[t]{0.20\textwidth}\includegraphics[width=\linewidth]{figure3b_static_lock.pdf}\caption{StaticLock}\end{subfigure}\hfill
\begin{subfigure}[t]{0.25\textwidth}\includegraphics[width=\linewidth]{figure3c_vcd_risk_coverage.pdf}\caption{VCD风险--覆盖率}\end{subfigure}\hfill
\begin{subfigure}[t]{0.32\textwidth}\includegraphics[width=\linewidth]{figure3d_temporal_strategies.pdf}\caption{时序合成}\end{subfigure}
\caption{组件归因。(a--b)中的连线连接同一动作片段；(c)显示各事件曲线及其中位数/四分位距包络；(d)采用对数误差轴展示全部片段。}
\label{fig:exp2}
\end{figure*}

采集时刻对齐将头动条件下的平移泄漏P95中位数由43.75~mm降至15.38~mm。StaticLock将中心化静止平移P95由13.73~mm降至0.82~mm，12个配对片段均呈同一方向。VCD的事件级曲线显示，高分观测集中于较低误差区域；AURC中位数为5.65~mm，而全覆盖风险为6.53~mm，12个事件均获得正向排序收益。

在时序合成比较中，延迟Linear/SLERP在360~ms有效滞后下获得8.93~mm的平移滞后对齐RMSE。平滑卡尔曼外推将滞后降至305~ms，但RMSE升至45.64~mm；Hermite插值与Linear/SLERP接近，未表现出稳定优势。该结果表明，较小的相位延迟并不自动带来更可靠的空间轨迹，历史插值在本系统的观测频率和运动范围内提供了更稳健的折中。

\subsection{实验三：感知质量与空间标注信任}

\paragraph{设计与任务。}
24名参与者完成$2\times3$完全被试内实验，比较One-Euro Anchor与\system在无线鼠标、订书机和游戏手柄上的表现。物体顺序、对象内方法顺序和匿名标签映射均衡。每个方法--物体区块包含三阶段：参与者首先在物体静止时通过左右转头和靠近观察锚点；随后拿起物体，完成约25--30~cm平移、约30--45$^{\circ}$旋转、短暂停留和重新放置；最后，实验员使用不透明挡板完全遮挡目标0.6--0.9~s，并让参与者观察恢复过程。

参与者中女性12名、男性12名，平均年龄为27.08岁（SD=4.45）；22名惯用右手，2名惯用左手。全部参与者完成实验并纳入分析。3名参与者报告实验后轻微不适，无人中止。两种方法的候选率分别为12.85和12.86~Hz，VCD接纳率为88.6\%和88.2\%，输出可用率为99.10\%和99.06\%。每种方法的72个遮挡区块均进入FrozenUncertain状态，未进入Lost。

\paragraph{测量与分析。}
七个研究定制条目分别测量静止稳定、位置正确、运动附着、姿态一致、恢复一致、依赖意愿和稳定--响应平衡。研究同时采用经物体锚定语境适配的Augmentation Quality嵌入质量与交互质量子量表、Trust in Automation的可靠性/能力与理解/可预测性分量表，以及S-TIAS信任量表。定制条目逐项分析；区块级结果先在参与者内跨三个物体取均值。两种方法使用双侧Wilcoxon符号秩检验，七项主条目与五项已发表量表分别进行Holm校正，并报告匹配秩双列相关$r_{rb}$。

\begin{figure*}[t]
\centering
\begin{subfigure}[t]{0.49\textwidth}\includegraphics[width=\linewidth]{figure4_exp3_primary_effects.pdf}\caption{主要锚点质量结局}\end{subfigure}\hfill
\begin{subfigure}[t]{0.49\textwidth}\includegraphics[width=\linewidth]{figure5_exp3_scale_effects.pdf}\caption{已发表量表结局}\end{subfigure}
\caption{实验三的配对差中位数与四分位距。正值表示\system评分更高；实心标记表示Holm校正后$p<.05$。}
\label{fig:exp3}
\end{figure*}

\begin{table*}[t]
\centering
\caption{实验三主要结果。配对差定义为EgoAnchor减One-Euro。}
\label{tab:exp3}
\small
\begin{tabularx}{\textwidth}{lccccl}
\toprule
结局 & One-Euro中位数 [Q1,Q3] & EgoAnchor中位数 [Q1,Q3] & 配对差中位数 & $p_{\mathrm{Holm}}$ & $r_{rb}$ \\
\midrule
静止稳定 & 4.33 [3.92,4.67] & 5.00 [4.58,5.67] & 1.00 & $<.001$ & .90 \\
位置正确 & 4.67 [4.00,5.00] & 5.00 [4.67,5.75] & 0.67 & .001 & .85 \\
运动附着 & 5.00 [4.33,5.08] & 5.17 [4.58,5.33] & 0.33 & .267 & .29 \\
姿态一致 & 4.67 [4.67,5.00] & 5.00 [4.67,5.33] & 0.33 & .162 & .41 \\
恢复一致 & 4.17 [3.92,4.33] & 5.00 [4.33,5.42] & 0.67 & $<.001$ & .95 \\
依赖意愿 & 4.00 [4.00,4.33] & 5.00 [4.33,5.33] & 1.00 & $<.001$ & .87 \\
稳定--响应平衡 & 4.50 [4.25,4.75] & 4.83 [4.58,5.33] & 0.33 & .003 & .90 \\
\midrule
AQ嵌入质量 & 4.78 [4.50,5.03] & 5.17 [4.50,5.47] & 0.44 & .004 & .75 \\
AQ交互质量 & 4.89 [4.64,5.14] & 5.06 [4.50,5.33] & 0.06 & .446 & .19 \\
TiA可靠性/能力 & 4.17 [3.83,4.50] & 4.58 [4.29,4.83] & 0.33 & $<.001$ & 1.00 \\
TiA理解/可预测性 & 4.00 [3.94,4.25] & 4.38 [4.00,5.00] & 0.38 & .004 & .75 \\
S-TIAS信任 & 5.33 [4.92,5.67] & 5.67 [5.33,6.00] & 0.50 & .004 & .79 \\
\bottomrule
\end{tabularx}
\end{table*}

\paragraph{结果。}
\system在静止稳定、位置正确、恢复一致、依赖意愿和稳定--响应平衡上获得更高评分（表~\ref{tab:exp3}）。其中，静止稳定、恢复一致和依赖意愿的配对差中位数分别为1.00、0.67和1.00，效应量均较大。运动附着和姿态一致未检出校正后的显著差异，表明用户侧优势主要集中在静止观察、放置后稳定和遮挡恢复，而非所有动态跟随维度。

在已发表量表中，\system提高了AQ嵌入质量、TiA可靠性/能力、TiA理解/可预测性和S-TIAS信任；AQ交互质量未检出差异。15名参与者总体偏好\system，4名偏好One-Euro，5名无明显偏好；18名更愿意信任\system的空间标注，1名更信任One-Euro，5名无明显差异。7名参与者的总体偏好与信任选择不一致，说明使用偏好与对空间信息的依赖意愿属于相关但不同的判断。

\section{讨论}

\subsection{动态锚定的核心是运行时语义}
实验结果支持将动态物体锚定定义为完整的运行时系统问题。感知后端回答“目标在某一图像时刻位于何处”，运行时进一步回答“该观测属于哪个世界时刻、是否应更新状态、渲染帧之间输出什么，以及追踪退化后应用应看到什么”。在实验二中保持候选流与参考轨迹不变，使采集时刻对齐、StaticLock和时序合成的作用能够被独立观察。由此，论文的主要贡献不在于重新设计单个视觉模型，而在于建立一套可由混合现实应用消费的观测到锚点契约。

\subsection{稳定性与响应性构成联合结果}
StaticLock和历史插值显著降低了静止泄漏、帧间抖动和滞后对齐残差，同时引入约数百毫秒的有效滞后与起动转换时间。这一工作点适合“观察--拿取--放置--读取”类任务：物体在操作前后存在明确静止阶段，放置后的稳定和遮挡恢复直接影响标注可读性。对于持续快速操控，当前时刻相位误差仍可能主导体验。后续系统可依据任务阶段在稳定优先和响应优先策略之间切换，而非为所有应用固定单一参数。

\subsection{可预测的恢复行为影响信任}
用户研究中最明显的提升出现在静止稳定、恢复一致和依赖意愿，而运动附着与AQ交互质量未出现显著差异。该模式表明，用户对空间标注的信任不仅取决于平均几何误差，也取决于锚点在静止和感知退化时是否表现出一致、可预期的行为。偏好与信任选择的分离进一步说明，单一总体偏好不足以描述用户是否愿意据此进行判断或操作。

\subsection{局限}
受控实验采用平台追踪控制器作为参考，因此结果反映同一Quest坐标体系内的相对系统行为，未覆盖平台参考与头显之间的共同漂移。用户研究评价了三类日常物体上的感知质量与信任，但未提供这些物体的外部绝对位姿真值或任务绩效。系统依赖刚性目标的三维模型，对严重遮挡、透明或镜面材质、低质量双目几何和模型误差仍较敏感。当前实现使用外部GPU工作站和无线网络，尚未形成完全头显端部署。部分适配量表的内部一致性随方法条件变化，完整$\alpha$与$\omega$结果应在补充材料中逐项报告。

\section{结论}
本文提出\system，一套将延迟、低频且质量不均的视觉6DoF观测转换为连续动态物体锚点的混合现实系统。系统通过采集时刻世界对齐恢复历史观测的世界坐标语义，以可靠性接纳限制有害更新，通过历史轨迹合成生成渲染帧级输出，并以StaticLock和生命周期管理定义静止、运动、遮挡和恢复阶段的行为。受控实验显示，\system显著改善静止世界一致性和滞后对齐后的动态轨迹质量；组件实验揭示了稳定性与响应性的具体来源；用户研究进一步表明，这些运行时特性能够提高用户对静止稳定、恢复一致和空间标注可信度的评价。动态物体锚定由此不再是离散位姿结果的简单显示，而是一项需要联合评价时空语义、连续性和可恢复性的系统能力。

% 待正式提交时补充匿名伦理声明与准确的AI辅助写作披露。
\bibliographystyle{abbrv-doi-hyperref-narrow}
\bibliography{egoanchor_cn_refs}
\end{document}
